\renewcommand{\thesection}{Appendix:}
\section[Lie algebra cohomology as a geometrical construction]{{}\hspace*{-5pt}Lie algebra cohomology as a geometrical construction}

Although in this article we shall need only the cohomology of an Abelian Lie
algebra, we take the opportunity here to describe the cohomology of a general
Lie algebra ${\mathfrak{g}}$ in terms of dif\/ferential geometry on~$G$, a Lie
group whose Lie algebra is~${\mathfrak{g}}$. We believe that for a general
parabolic geometry, we shall need this geometric interpretation for a nilpotent
Lie algebra. Suppose ${\mathbb{V}}$ is a $G$-module and use the same notation
for the corresponding representation of~${\mathfrak{g}}$. Following but
adapting~\cite{ce}, we are going to present the Lie algebra cohomology
$H^r({\mathfrak{g}},{\mathbb{V}})$ as a geometrical construction on~$G$. Beware
that $G$ is no longer the Lie group ${\mathrm{SL}}(n+1,{\mathbb{R}})$ as it was
until now. This section is written to be self-contained with the aim of being
useful elsewhere. This material is well-known to experts and implicit
in~\cite{ce} but we believe it worthwhile laying out the details.

We shall view $G$ as a homogeneous space under its own action on the left. Its
tangent bundle~$TG$ is then regarded as a homogeneous bundle and can be
identif\/ied as $G\times{\mathfrak{g}}$, where ${\mathfrak{g}}$ is the Lie
algebra of $G$. It is convenient to write this isomorphism as
\begin{equation}\label{m-c}
X\mapsto X\intprod\theta \qquad\mbox{for vector f\/ields $X$ on $G$},
\end{equation}
where $\theta$ is a $1$-form on $G$ with values in~${\mathfrak{g}}$ known as
the {\em Maurer--Cartan\/} form~\cite{ss}. To compute with $\theta$ it is
convenient to write functions on $G$ with values in ${\mathfrak{g}}$ as
$X^\alpha$ and then (\ref{m-c}) becomes
\[
X^a\mapsto X^\alpha\equiv\theta_a^\alpha X^a\qquad \mbox{with inverse}
\quad X^\alpha\mapsto X^a\equiv\phi_\alpha^a X^\alpha,
\]
where $\phi_\alpha^a\theta^\beta_a=\delta_\alpha{}^\beta$ and $\theta_a^\alpha\phi_\alpha^b=\delta_a{}^b$.
A vector f\/ield $X^a$ on $G$ is {\em left-invariant\/} if and only if the
corresponding function $X^\alpha:G\to{\mathfrak{g}}$ is constant. Choosing any
torsion-free af\/f\/ine connection $\nabla_a$ on $G$ and expanding the def\/inition
$\nabla_a(\theta_b^\alpha X^b)=0$, the left-invariant vector f\/ields are those
that satisfy
\[\theta_b^\alpha \nabla_aX^b-X^b \nabla_a\theta_b^\alpha=0\]
from which it follows easily that the Lie bracket of two left-invariant
vector f\/ields is again left-invariant.
Since the left-invariant vector f\/ields on $G$ are of the form $\phi(X)$ for
$X\in{\mathfrak{g}}$ we may def\/ine the Lie bracket
on ${\mathfrak{g}}$ by transportation: %--
\begin{equation}\label{defnLie}
\phi([X,Y])=[\phi(X),\phi(Y)]\qquad\mbox{for} \quad X,Y\in{\mathfrak{g}}.
\end{equation}
For computational purposes, let us write
$[X,Y]^\gamma=\Gamma_{\alpha\beta}{}^\gamma X^\alpha Y^\beta$ for the Lie
bracket on~${\mathfrak{g}}$. Then we can write out~(\ref{defnLie}) explicitly
as
 \begin{gather*}
\phi_\gamma^c\Gamma_{\alpha\beta}{}^\gamma X^\alpha Y^\beta =
[\phi_\alpha^a X^\alpha,\phi_\beta^b Y^\beta]^c
 = \phi_\alpha^a X^\alpha \nabla_a(\phi_\beta^c Y^\beta)-
\phi_\alpha^b Y^\alpha \nabla_b(\phi_\beta^c X^\beta)\\
\hphantom{\phi_\gamma^c\Gamma_{\alpha\beta}{}^\gamma X^\alpha Y^\beta =
[\phi_\alpha^a X^\alpha,\phi_\beta^b Y^\beta]^c}{}
 = (\phi_\alpha^a \nabla_a\phi_\beta^c-\phi_\beta^a \nabla_a\phi_\alpha^c)
X^\alpha Y^\beta
\end{gather*}
or, in other words, as
\[
\Gamma_{\alpha\beta}{}^\gamma\phi_\gamma^c
=\phi_\alpha^a \nabla_a\phi_\beta^c-\phi_\beta^a \nabla_a\phi_\alpha^c.
\]
Bearing in mind that $\phi_\beta^b\theta_b^\gamma=\delta_\beta{}^\gamma$ whence
$\theta_b^\gamma \nabla_a\phi_\beta^b+\phi_\beta^b \nabla_a\theta_b^\gamma=0$,
we may rewrite this as
\[\Gamma_{\alpha\beta}{}^\gamma
=-\phi_\alpha^a\phi_\beta^b \nabla_a\theta_b^\gamma+
\phi_\beta^a\phi_\alpha^b \nabla_a\theta_b^\gamma
=-\phi_\alpha^a\phi_\beta^b(\nabla_a\theta_b^\gamma-\nabla_b\theta_a^\gamma)\]
and f\/inally as
\begin{equation}\label{m-c-eq1}
\nabla_a\theta_b^\gamma-\nabla_b\theta_a^\gamma
+\Gamma_{\alpha\beta}{}^\gamma\theta_a^\alpha\theta_b^\beta=0.
\end{equation}
This formula employs an arbitrary torsion-free connection. Without indices and
without this connection  it is more usually written as
\begin{equation}\label{m-c-eq2}  d\theta+\tfrac12[\theta,\theta]=0\qquad
\mbox{or}\qquad d\theta+\theta\wedge\theta=0.
\end{equation}
In other words, the def\/inition (\ref{defnLie}) of the Lie bracket on
${\mathfrak{g}}$ is equivalent to (\ref{m-c-eq1}) or~(\ref{m-c-eq2}), usually
known as the {\em Maurer--Cartan\/} equation~\cite{ss}.

Now suppose ${\mathbb{V}}$ is a $G$-module and use the same notation for the
corresponding ${\mathfrak{g}}$-module. There are two canonically def\/ined
connections on the vector bundle $V=G\times{\mathbb{V}}$ over~$G$. One of them
is the evident f\/lat connection, ignoring the action of $G$ on ${\mathbb{V}}$.
We shall denote it by $d$ since it is the exterior derivative acting on
functions with values in~${\mathbb{V}}$. The other one takes the isomorphism
from (\ref{twist})
\begin{equation}\label{twistagain}
V=G\times{\mathbb{V}}\cong G\times{\mathbb{V}}\qquad\mbox{by}\quad
(g,v)\mapsto(g,gv)\end{equation}
and pulls back the evident f\/lat connection on the right hand side as was done
in Section~\ref{outline} and we shall denote this one by $\nabla$ as was done there.
To relate these two connections more explicitly suppose $f:G\to{\mathbb{V}}$ is
a section of $V$ that is constant after the twisting~(\ref{twistagain}). It
means that the function $g\mapsto gf(g)$ is constant. If so, then for f\/ixed
$g\in G$ and~$X\in{\mathfrak{g}}$, the function
\[{\mathbb{R}}\ni t\mapsto ge^{tX}f\big(ge^{tX}\big)\in{\mathbb{V}}\]
is constant. Equivalently, the function $t\mapsto e^{tX}f(ge^{tX})$ is
constant and so
\[0=\frac{d}{dt}\big(e^{tX}f\big(ge^{tX}\big)\big)\Big|_{t=0}
=\frac{d}{dt}f\big(ge^{tX}\big)\big|_{t=0}+Xf(g)=(\phi(X)f+Xf)(g),\]
where this last equality is due to the f\/low of the left-invariant vector f\/ield
$\phi(X)$ being the one-parameter subgroup of right-translations
$g\mapsto ge^{tX}$ (see, e.g.~\cite{w}). For computational purposes, let us
write $X^\alpha\mapsto X^\alpha\rho_\alpha$ where
$\rho_\alpha\in{\mathfrak{g}}^*\otimes\End({\mathbb{V}})$ for the action of
${\mathfrak{g}}$ on~${\mathbb{V}}$. Then $g\mapsto gf(g)$ is constant if and
only if
\[0=\phi(X)f+Xf=X^\alpha\phi_\alpha^ad_af+X^\alpha\rho_\alpha f
=X^a(d_af+\theta_a^\alpha\rho_\alpha f)\]
for all left-invariant vector f\/ields~$X^a$. It follows that the connection
$\nabla_a$ on $V$ is given by
\begin{equation}\label{nabla}
f\mapsto\nabla_af=d_af+\theta_a^\alpha\rho_\alpha f\qquad
\mbox{or, without indices, as}\quad
f\mapsto\nabla f=df+\theta f,\end{equation}
where $\theta\in\Lambda^1\otimes{\mathfrak{g}}$ is the Maurer--Cartan form. As a
check, the dif\/ferential in the coupled de~Rham sequence (\ref{coupled}) is
\[\Lambda^p\otimes{\mathbb{V}}\ni\omega\stackrel{\nabla}{\longmapsto}
d\omega+\theta\wedge\omega
\in\Lambda^{p+1}\otimes{\mathbb{V}}\]
and the Maurer--Cartan equation~(\ref{m-c-eq2}) shows that the composition
${\mathbb{V}}\stackrel{\nabla}{\to}\Lambda^1\otimes{\mathbb{V}}
\stackrel{\nabla}{\to}\Lambda^2\otimes{\mathbb{V}}$ is given by
\[\nabla^2f=d(df+\theta f)+\theta\wedge(df+\theta f)
=(d\theta+\theta\wedge\theta)f=0\]
and the connection $\nabla$ is f\/lat, as expected. More generally, the whole
sequence
\begin{equation}\label{coupledtobbV}
0\to{\mathbb{V}}\stackrel{\nabla}{\longrightarrow}
\Lambda^1\otimes{\mathbb{V}}\stackrel{\nabla}{\longrightarrow}
\Lambda^2\otimes{\mathbb{V}}
\stackrel{\nabla}{\longrightarrow}\cdots\stackrel{\nabla}{\longrightarrow}
\Lambda^p\otimes{\mathbb{V}}
\stackrel{\nabla}{\longrightarrow}
\Lambda^{p+1}\otimes{\mathbb{V}}
\stackrel{\nabla}{\longrightarrow}
\cdots\end{equation}
is a complex.

Consider the linear mapping
\[\Lambda^p{\mathfrak{g}}^*\otimes{\mathbb{V}}\ni v_{\alpha\beta\cdots\gamma}
\stackrel{\theta}{\longmapsto}v_{ab\cdots c}\equiv
\theta_a^\alpha\theta_b^\beta\cdots\theta_c^\gamma
v_{\alpha\beta\cdots\gamma}\in\Gamma(G,\Lambda^p\otimes{\mathbb{V}})\]
in which ${\mathbb{V}}$ is just a passenger (i.e.~plays no r\^{o}le). We shall
refer to the resulting ${\mathbb{V}}$-valued $p$-form as {\em left-invariant\/}
just as we would if ${\mathbb{V}}$ were absent.
\begin{lemma} The connection
$\nabla:\Lambda^p\otimes{\mathbb{V}}\to\Lambda^{p+1}\otimes{\mathbb{V}}$
preserves left-invariance.\end{lemma}
\begin{proof}We use the Maurer--Cartan equation (\ref{m-c-eq1}) to compute
\begin{gather}
\nabla_{[a}v_{bc\cdots d]} =d_{[a}(\theta_b^\beta\theta_c^\gamma\cdots\theta_{d]}^\delta
v_{\beta\gamma\cdots\delta})+
\theta_{[a}^\alpha\theta_b^\beta\theta_c^\gamma\cdots\theta_{d]}^\delta
\rho_\alpha v_{\beta\gamma\cdots\delta}
\nonumber\\
 \hphantom{\nabla_{[a}v_{bc\cdots d]}}{} = p(d_{[a}\theta_b^\epsilon)\theta_c^\gamma\cdots\theta_{d]}^\delta
v_{\epsilon\gamma\cdots\delta}+
\theta_{[a}^\alpha\theta_b^\beta\theta_c^\gamma\cdots\theta_{d]}^\delta
\rho_\alpha v_{\beta\gamma\cdots\delta}
\nonumber\\
 \hphantom{\nabla_{[a}v_{bc\cdots d]}}{} =-(p/2)\Gamma_{\alpha\beta}{}^\epsilon
\theta_{[a}^\alpha\theta_b^\beta\theta_c^\gamma\cdots\theta_{d]}^\delta
v_{\epsilon\gamma\cdots\delta}+
\theta_{[a}^\alpha\theta_b^\beta\theta_c^\gamma\cdots\theta_{d]}^\delta
\rho_\alpha v_{\beta\gamma\cdots\delta}
\nonumber\\
 \hphantom{\nabla_{[a}v_{bc\cdots d]}}{} =\theta_{[a}^\alpha\theta_b^\beta\theta_c^\gamma\cdots\theta_{d]}^\delta
\rho_\alpha v_{\beta\gamma\cdots\delta}
+(-1)^p(p/2)
\theta_{[a}^\alpha\theta_b^\beta\theta_c^\gamma\cdots\theta_{d]}^\delta
\Gamma_{\alpha\beta}{}^\epsilon v_{\gamma\cdots\delta\epsilon}
\nonumber\\
 \hphantom{\nabla_{[a}v_{bc\cdots d]}}{} = \theta_a^\alpha\theta_b^\beta\theta_c^\gamma\cdots\theta_d^\delta
\left(\rho_{[\alpha}v_{\beta\gamma\cdots\delta]}
+(-1)^p(p/2)
\Gamma_{[\alpha\beta}{}^\epsilon v_{\gamma\cdots\delta]\epsilon}\right),\label{key}
\end{gather}
as required.
\end{proof}

In fact (\ref{key}) shows that $\nabla\theta v=\theta\partial v$, where
\begin{gather}
\partial: \ \Lambda^p{\mathfrak{g}}^*\otimes{\mathbb{V}}\to
\Lambda^{p+1}{\mathfrak{g}}^*\otimes{\mathbb{V}}
\qquad\mbox{is given by}\nonumber\\
\qquad
v_{\beta\gamma\cdots\delta}\mapsto\rho_{[\alpha}v_{\beta\gamma\cdots\delta]}
+(-1)^p(p/2)\Gamma_{[\alpha\beta}{}^\epsilon v_{\gamma\cdots\delta]\epsilon}.\label{thisispartial}
\end{gather}
It also follows that
\begin{equation}\label{Koszul}
0\to{\mathbb{V}}\stackrel{\partial}{\longrightarrow}
{\mathfrak{g}}^*\otimes{\mathbb{V}}\stackrel{\partial}{\longrightarrow}
\Lambda^2{\mathfrak{g}}^*\otimes{\mathbb{V}}
\stackrel{\partial}{\longrightarrow}\cdots\stackrel{\partial}{\longrightarrow}
\Lambda^p{\mathfrak{g}}^*\otimes{\mathbb{V}}
\stackrel{\partial}{\longrightarrow}
\Lambda^{p+1}{\mathfrak{g}}^*\otimes{\mathbb{V}}
\stackrel{\partial}{\longrightarrow}
\cdots\end{equation}
is a complex of~${\mathfrak{g}}$-modules. Alternatively, this may be directly
verif\/ied from (\ref{thisispartial}) using
\begin{itemize}\itemsep=0pt
\item
$\rho_{[\alpha}\rho_{\beta]}=\frac12\Gamma_{\alpha\beta}{}^\gamma\rho_\gamma$
(i.e.~that $\rho:{\mathfrak{g}}\to\End({\mathbb{V}})$ is a representation),
\item $\Gamma_{[\alpha\beta}{}^\delta\Gamma_{\gamma]\delta}{}^\epsilon=0$
(i.e.~the Jacobi identity in ${\mathfrak{g}}$),
\item $(Xv)_{\beta\gamma\cdots\delta}
=X^\alpha\rho_\alpha v_{\beta\gamma\cdots\delta}
+(-1)^ppX^\alpha\Gamma_{\alpha[\beta}{}^\epsilon
v_{\gamma\cdots\delta]\epsilon}$ (the action of ${\mathfrak{g}}$ on
$\Lambda^p{\mathfrak{g}}^*\otimes{\mathbb{V}}$).
\end{itemize}
We have shown that there are two ways of def\/ining Lie algebra cohomology as
follows.
\begin{theorem}\label{geometricrealisation} The Lie algebra cohomology
$H^r({\mathfrak{g}},{\mathbb{V}})$ may be defined as either
\begin{itemize}\itemsep=0pt
\item the cohomology of the Koszul complex~\eqref{Koszul}, or
\item the cohomology of the complex \eqref{coupledtobbV} restricted to
left-invariant forms.
\end{itemize}
\end{theorem}
\begin{remark} Although we use this theorem in the main body of this article,
it is easily avoided. In tackling a general parabolic geometry,
however, we believe that Theorem~\ref{geometricrealisation} will be essential.
\end{remark}

\begin{remark}
As a minor variation on this construction, suppose $G$ is enlarged to~$Q$, a
semi-direct product
\[Q=G_0\ltimes G\qquad\mbox{i.e.}\qquad
\begin{picture}(0,0)
\put(0,0){
${\mathrm{Id}}\to G\lhd Q\longrightarrow G_0\to{\mathrm{Id}}$}
\qbezier (67,-0.5) (74,-4) (81,-0.5)
\put(65,0){\vector(-3,1){0}}
\end{picture}\]
and suppose that ${\mathbb{V}}$ extends to a representation of~$Q$. We identify
$G$ with the $Q$-homogeneous space~$Q/G_0$, noting that when the action of $Q$
on $G=Q/G_0$ is restricted to~$G$ it coincides with its usual action of $G$ on
itself by left translation. The $Q$-homogeneous bundle
$V\equiv Q\times_{G_0}{\mathbb{V}}$ on~$Q/G_0$ is equipped with a f\/lat
connection $\nabla$ by dint of the canonical trivialisation
\[V=Q\times_{G_0}{\mathbb{V}}\ni(q,v)\mapsto(qG_0,qv)\in
Q/G_0\otimes{\mathbb{V}}=G\times{\mathbb{V}},\]
which clearly coincides with $\nabla$ def\/ined by~(\ref{twistagain}). This
connection is $Q$-equivariant. Consequently, not only does the twisted
de~Rham complex
\[V\xrightarrow{\,\nabla\,}\Lambda^1\otimes V
\xrightarrow{\,\nabla\,}\Lambda^2\otimes V
\xrightarrow{\,\nabla\,}\cdots
\xrightarrow{\,\nabla\,}\Lambda^p\otimes V
\xrightarrow{\,\nabla\,}\Lambda^{p+1}\otimes V
\xrightarrow{\,\nabla\,}\cdots\]
coincide with (\ref{coupledtobbV}) and thereby compute the Lie algebra
cohomology $H^r({\mathfrak{g}},{\mathbb{V}})$ when restricted to $G$-invariant
forms, but also the complex (\ref{Koszul}) is automatically one of
$Q$-modules where the $Q$-action on ${\mathbb{V}}$
is as supposed and the $Q$-action on ${\mathfrak{g}}^*$ is induced by the
conjugation action of $Q$ on $G$ (bearing in mind that $G$ is a normal
subgroup of~$Q$).
\end{remark}

\renewcommand{\thesubsection}{Addendum:}

\subsection[A canonical connection on $G$]{{}\hspace*{-3pt}A canonical connection on $\boldsymbol{G}$}

Again, although it is unnecessary for the current article and already known to
experts, we suspect that the following optional extra will be invaluable in
dealing with a general parabolic geometry. Since we already have established
suitable notation in the Appendix above, we take the opportunity of presenting
it here. Our canonical connection $D_a$ was introduced in~\cite{cs} as the
`$(+)$-connection' and Lemma~\ref{torsion_is_bracket} is stated without
proof as~\cite[Proposition~2.12]{knII}.

The trivialisation $T^*G=G\times{\mathfrak{g}}^*$ provided by the Maurer--Cartan
form also equips $G$ with a~canonical f\/lat af\/f\/ine connection $D_a$ def\/ined
by
\[D_a\omega_b\equiv\theta_b^\beta d_a\omega_\beta,\]
where $\omega_\beta\equiv\phi_\beta^c\omega_c$ and
$d_a$ on the right hand side of this equation simply takes the
gradient of a~function with values in~${\mathfrak{g}}^*$. If we expand using
any torsion-free af\/f\/ine connection $\nabla_a$
\[D_a\omega_b=\theta_b^\beta \nabla_a(\phi_\beta^c\omega_c)
=\nabla_a\omega_b+(\theta_b^\beta \nabla_a\phi_\beta^c)\omega_c
=\nabla_a\omega_b-(\phi_\beta^c\nabla_a\theta_b^\beta)\omega_c,\]
then we see that, for $f$ a smooth function,
\[D_aD_bf-D_bD_af
=(-\phi_\beta^c \nabla_a\theta_b^\beta+\phi_\beta^c \nabla_b\theta_a^\beta)D_cf
=-\phi_\gamma^c(\nabla_a\theta_b^\gamma-\nabla_b\theta_a^\gamma)D_cf\]
and so the canonical connection $D_a$ has torsion
\[T_{ab}{}^c=-(\nabla_a\theta_b^\gamma-\nabla_b\theta_a^\gamma)\phi_\gamma^c.\]
Alternatively, from (\ref{m-c-eq1}) we see that
\[T_{ab}{}^c\theta_c^\gamma=-(\nabla_a\theta_b^\gamma-\nabla_b\theta_a^\gamma)
=\Gamma_{\alpha\beta}{}^\gamma\theta_a^\alpha\theta_b^\beta,\]
which we record as the following lemma.
\begin{lemma}\label{torsion_is_bracket}
The torsion of\/ $D_a$ coincides with
the Lie bracket on\/~${\mathfrak{g}}$ under the Maurer--Cartan parallelism.
\end{lemma}
Notice that $D_a\theta_b^\beta=0$. It is another way to
characterise~$D_a$ and, indeed, is the main point of this construction as
follows.
\begin{lemma} Even locally, the kernel of the induced operator
\[D: \ \Lambda^p\to\Lambda^1\otimes\Lambda^p\]
is the left-invariant forms on~$G$.
\end{lemma}
\begin{proof}
Recall that the left-invariant forms are obtained as
\[\Lambda^p{\mathfrak{g}}^*\ni v_{\alpha\beta\cdots\gamma}
\stackrel{\theta}{\longmapsto}v_{ab\cdots c}\equiv
\theta_a^\alpha\theta_b^\beta\cdots\theta_c^\gamma
v_{\alpha\beta\cdots\gamma}\in\Gamma(G,\Lambda^p)\]
and it clear that such forms are annihilated by~$D$. Conversely, since $D$ is
f\/lat and all covariant constant sections are already accounted for, there can
be no more, even locally.
\end{proof}
\begin{remark} Of course, this lemma also holds for ${\mathbb{V}}$-valued
dif\/ferential forms where the connection is trivially coupled
with~${\mathbb{V}}$ and it is this that we have in mind in constructing
the BGG complex in general.
\end{remark}
